\chapter*{ภาคผนวก ก\quad คู่มือการคำนวณความคลาดเคลื่อน การเขียนกราฟ และเกณฑ์รูบริกส์}
\addcontentsline{toc}{chapter}{ภาคผนวก ก\quad คู่มือการคำนวณความคลาดเคลื่อน การเขียนกราฟ และเกณฑ์รูบริกส์}
\markboth{ภาคผนวก ก\quad คู่มือการคำนวณความคลาดเคลื่อนและการเขียนกราฟ}{ภาคผนวก ก\quad คู่มือการคำนวณความคลาดเคลื่อนและการเขียนกราฟ}

\section*{1. สูตรการหาค่าเฉลี่ยและเปอร์เซ็นต์ความคลาดเคลื่อนสำหรับรายงานปฏิบัติการ}

ในการทดลองทางฟิสิกส์และจุลชีววิทยา เราจะทำการวัดซ้ำอย่างน้อย $3$ ครั้ง ($x_1, x_2, x_3$) เพื่อลดความคลาดเคลื่อนแบบสุ่ม โดยมีขั้นตอนการคำนวณ 2 ขั้นตอนง่ายๆ ดังนี้
\begin{enumerate}
    \item \textbf{การหาค่าเฉลี่ยเลขคณิต} ($\bar{x}$)
    \begin{equation*}
        \bar{x} = \frac{x_1 + x_2 + x_3}{3}
    \end{equation*}
    \item \textbf{การหาเปอร์เซ็นต์ความคลาดเคลื่อนสัมพัทธ์} (เทียบกับค่ามาตรฐาน $x_{\text{std}}$)
    \begin{equation*}
        \text{เปอร์เซ็นต์ความคลาดเคลื่อน} = \frac{|\bar{x} - x_{\text{std}}|}{x_{\text{std}}} \times 100\%
    \end{equation*}
\end{enumerate}

\section*{2. หลักการเขียนกราฟทางวิทยาศาสตร์และการหาความชัน (Slope)}

เมื่อต้องเขียนกราฟความสัมพันธ์ระหว่างตัวแปร 2 ตัว (เช่น กราฟระหว่างระยะทาง $s$ กับเวลา $t$ ในปฏิบัติการที่ 2 หรือกราฟมาตรฐานระหว่างค่า $\text{OD}_{600}$ กับความเข้มข้นเซลล์ $c$ ในปฏิบัติการที่ 12) ให้ปฏิบัติตามหลักการ 4 ข้อต่อไปนี้
\begin{enumerate}
    \item \textbf{การเลือกแกน}\quad ให้วาง\textbf{ตัวแปรต้น} (ตัวแปรที่เราเป็นผู้กำหนดค่า เช่น เวลา $t$ หรือความเข้มข้น $c$) ไว้บน\textbf{แกนนอน (แกน X)} และวาง\textbf{ตัวแปรตาม} (ค่าที่อ่านได้จากเครื่องมือวัด เช่น ระยะทาง $s$ หรือค่า $\text{OD}_{600}$) ไว้บน\textbf{แกนตั้ง (แกน Y)}
    \item \textbf{การระบุชื่อแกนและหน่วย}\quad ต้องเขียนชื่อปริมาณและระบุหน่วยในวงเล็บกำกับที่ปลายแกนทั้งสองเสมอ เช่น $\text{เวลา } t\text{ (s)}$ และ $\text{ระยะทาง } s\text{ (cm)}$
    \item \textbf{การลากเส้นแนวโน้ม (Best-Fit Line)}\quad ห้ามลากเส้นต่อจุดแบบฟันปลาไปมา แต่ให้ใช้ไม้บรรทัดทาบแล้วลาก\textbf{เส้นตรงเพียงเส้นเดียว}ให้ผ่านใกล้จุดข้อมูลทั้งหมดอย่างสมดุลที่สุด
    \item \textbf{การคำนวณความชันของกราฟเส้นตรง (Slope)}\quad เลือกจุด 2 จุดบนเส้นแนวโน้ม คือ $(x_1, y_1)$ และ $(x_2, y_2)$ แล้วคำนวณจากอัตราส่วน
    \begin{equation*}
        \text{ความชัน (Slope)} = \frac{\Delta y}{\Delta x} = \frac{y_2 - y_1}{x_2 - x_1}
    \end{equation*}
\end{enumerate}

\section*{3. เกณฑ์รูบริกส์การประเมินทักษะปฏิบัติการและรายงาน (CLO2, CLO3, CLO4)}

\begin{table}[H]
\centering
\caption{เกณฑ์การให้คะแนนปฏิบัติการรายสัปดาห์ตามผลลัพธ์การเรียนรู้ (CLO2, CLO3 และ CLO4)}
\footnotesize\setlength{\tabcolsep}{4pt}
\begin{tabularx}{\textwidth}{>{\raggedright\arraybackslash\bfseries}p{2.6cm} >{\raggedright\arraybackslash}X >{\raggedright\arraybackslash}X >{\raggedright\arraybackslash}X}
\toprule
ด้านที่ประเมิน & ระดับดีเยี่ยม (เต็มเกณฑ์) & ระดับดี (ผ่านเกณฑ์มาตรฐาน) & ระดับต้องปรับปรุง \\
\midrule
ทักษะการใช้เครื่องมือและการแก้ปัญหา (CLO2) & ตั้งค่าและใช้งานเครื่องมือ (เช่น Micropipette, Centrifuge, กล้องจุลทรรศน์ 100x, Spectrophotometer) ได้ถูกต้องตามหลักฟิสิกส์ทุกขั้นตอน & ใช้งานเครื่องมือได้ถูกต้อง แต่ยังต้องได้รับคำแนะนำในการปรับตั้งค่าหรือการแปลงหน่วยบางจุด & ใช้งานเครื่องมือผิดขั้นตอน เช่น ไม่ถ่วงสมดุลโรเตอร์ หรือไม่ทำ Blank ก่อนวัดค่า $\text{OD}_{600}$ \\
\addlinespace[4pt]
ความปลอดภัยและจรรยาบรรณข้อมูล (CLO3) & สวมอุปกรณ์ป้องกันครบถ้วน ปฏิบัติตามกฎความปลอดภัย 100\% บันทึกตัวเลขจากการวัดจริงด้วยความซื่อสัตย์และดูแลทำความสะอาดเครื่องมือหลังใช้ & ปฏิบัติตามกฎความปลอดภัยและบันทึกข้อมูลตามจริง แต่ทำความสะอาดหรือเก็บอุปกรณ์ยังไม่ครบทุกชิ้น & ประมาทในการใช้อุปกรณ์ความร้อน/ไฟฟ้า หรือมีการดัดแปลงตัวเลขผลการทดลองให้ตรงทฤษฎี \\
\addlinespace[4pt]
การทำงานเป็นทีมและรายงานปฏิบัติการ (CLO4) & สมาชิกทุกคนมีส่วนร่วมในการทดลอง ตารางบันทึกผลระบุหน่วยครบ กราฟถูกต้อง และอภิปรายผลเชื่อมโยงสู่งานจุลชีววิทยาได้ชัดเจน & สมาชิกช่วยกันทำงานดี รายงานครบทุกหัวข้อ แต่การอภิปรายผลความคลาดเคลื่อนยังสั้นเกินไป & สมาชิกบางคนไม่ร่วมทำการทดลอง หรือส่งรายงานล่าช้า/ขาดกราฟและการสรุปผล \\
\bottomrule
\end{tabularx}
\end{table}
